\chapter{Riemannian Metrics}
\label{chap:dc-riemannian-metrics}

Throughout this chapter, manifolds are Hausdorff and have a countable basis;
``differentiable'' means $C^\infty$, and $M^n$ denotes an $n$-dimensional
manifold $M$.

\section{Riemannian Metrics}

\begin{definition}[Riemannian metric]
  \label{def:dc-ch1-2-1}
  \lean{Riemannian.RiemannianMetric}
  \leanok
  \dcref{ch1:2.1}
  A \emph{Riemannian metric} (or \emph{Riemannian structure}) on a differentiable
  manifold $M$ assigns an inner product $\langle\ ,\ \rangle_p$ on each $T_pM$,
  varying differentiably with $p$. In a coordinate system
  $\mathbf{x}:U\subset\mathbb{R}^n\to M$, for
  $q=\mathbf{x}(x_1,\ldots,x_n)$ and
  $\frac{\partial}{\partial x_i}(q)=d\mathbf{x}_{(x_1,\ldots,x_n)}e_i$,
  the coefficients
  $$
  g_{ij}(x_1,\ldots,x_n)=
  \left\langle\frac{\partial}{\partial x_i}(q),
  \frac{\partial}{\partial x_j}(q)\right\rangle_q
  $$
  are differentiable on $U$. This condition is independent of the chart;
  equivalently, $\langle X,Y\rangle$ is differentiable for differentiable vector
  fields $X,Y$. The functions $g_{ij}=g_{ji}$ are the \emph{local representation
  of the metric}; $M$ with such a metric is a \emph{Riemannian manifold}.
\end{definition}

\begin{definition}[isometry]
  \label{def:dc-ch1-2-2}
  \lean{Riemannian.DCIsometry, Riemannian.DCIsometry.preservesMetric}
  \uses{def:dc-ch1-2-1}
  \leanok
  \dcref{ch1:2.2}
  For Riemannian manifolds $M,N$, a diffeomorphism $f:M\to N$ is an
  \emph{isometry} when
  $$
  \langle u,v\rangle_p=\langle df_p(u),df_p(v)\rangle_{f(p)},\qquad\text{(1)}
  $$
  for every $p\in M$ and $u,v\in T_pM$.
\end{definition}

\begin{definition}[local isometry]
  \label{def:dc-ch1-2-3}
  \lean{Riemannian.DCIsLocalIsometryAtFull, Riemannian.DCIsLocalIsometryAtFull.isLocalDiffeomorphAt, Riemannian.DCIsLocalIsometryAtFull.preservesMetricNear}
  \uses{def:dc-ch1-2-1, def:dc-ch1-2-2}
  \leanok
  \dcref{ch1:2.3}
  A differentiable map $f:M\to N$ between Riemannian manifolds is a
  \emph{local isometry at} $p\in M$ if some neighborhood $U$ of $p$ is mapped
  diffeomorphically onto $f(U)$ and the restriction satisfies (1). The manifold
  $M$ is \emph{locally isometric} to $N$ if every $p\in M$ has such a neighborhood.
\end{definition}

\begin{example}[Euclidean space]
  \label{ex:dc-ch1-2-4}
  \lean{Riemannian.DCEuclideanMetric}
  \uses{def:dc-ch1-2-1}
  \leanok
  \dcref{ch1:2.4}
  On $\mathbb{R}^n$, identify $\frac{\partial}{\partial x_i}$ with the standard
  basis vector $e_i$ and set $\langle e_i,e_j\rangle=\delta_{ij}$. This is
  Euclidean space of dimension $n$ with its standard Riemannian metric.
\end{example}

\begin{example}[immersed manifolds]
  \label{ex:dc-ch1-2-5}
  \lean{Riemannian.DCInducedForm, Riemannian.DCInducedMetric}
  \uses{def:dc-ch1-2-1, ex:dc-ch1-2-4}
  \leanok
  \dcref{ch1:2.5}
  If $f:M^n\to N^{n+k}$ is an immersion into a Riemannian manifold, define
  $$
  \langle u,v\rangle_p=\langle df_p(u),df_p(v)\rangle_{f(p)}.
  $$
  Injectivity of $df_p$ gives positive definiteness, and differentiability follows
  in coordinates; this is the metric \emph{induced by $f$}, and $f$ is an
  \emph{isometric immersion}. If $h:M^{n+k}\to N^k$ has regular value $q$, then
  $h^{-1}(q)$ is an $n$-submanifold with the metric induced by inclusion. For
  $h:\mathbb{R}^n\to\mathbb{R}$, $h(x)=\sum_i x_i^2-1$, the regular level
  $h^{-1}(0)=S^{n-1}$ has its canonical metric.
\end{example}
\begin{proof}
  \leanok
  \uses{ex:dc-ch1-2-4}
  \sketch
  Pull back the target inner product along $df_p$. Bilinearity and symmetry are
  inherited. If $u\ne0$, injectivity gives $df_p(u)\ne0$, hence
  $\langle u,u\rangle_p>0$. In coordinates the resulting coefficients are obtained
  by composing the differentiable target coefficients with the coordinate
  differential of $f$, so they vary differentiably and satisfy
  \cref{def:dc-ch1-2-1}.
\end{proof}

\begin{example}[Lie groups]
  \label{ex:dc-ch1-2-6}
  \uses{prop:dc-ch0-5-4}
  \notready
  \dcref{ch1:2.6}
  \notready
  A \emph{Lie group} is a group $G$ with a differentiable structure for which
  $(x,y)\mapsto xy^{-1}$ is differentiable. Left and right translations
  $L_x(y)=xy$ and $R_x(y)=yx$ are diffeomorphisms. A metric is \emph{left invariant}
  if every $L_x$ is an isometry, \emph{right invariant} if every $R_x$ is an
  isometry, and \emph{bi-invariant} if both conditions hold.

  A vector field $X$ is left invariant when $dL_xX=X$ for every $x$; it is determined
  by $X_e$. The bracket of left-invariant fields is left invariant, so
  $[X_e,Y_e]=[X,Y]_e$ defines the Lie algebra $\mathcal{G}=T_eG$. Given an inner
  product on $\mathcal{G}$, the formula
  $$
  \langle u,v\rangle_x=
  \left\langle(dL_{x^{-1}})_x(u),(dL_{x^{-1}})_x(v)\right\rangle_e
  \qquad\text{(2)}
  $$
  defines a left-invariant metric; the right-invariant construction is analogous,
  and every compact Lie group has a bi-invariant metric. For a bi-invariant metric,
  $$
  \langle[U,X],V\rangle=-\langle U,[V,X]\rangle
  \qquad\text{for }U,V,X\in\mathcal{G}.\tag{3}
  $$
  Conversely, a positive bilinear form on $\mathcal{G}$ satisfying (3), extended
  by (2), gives a bi-invariant metric.
\end{example}
\begin{proof}
  \sketch
  For left-invariant $X,Y$,
  $$
  dL_x[X,Y]f=[X,Y](f\circ L_x)=X(dL_xY)f-Y(dL_xX)f=[X,Y]f,
  $$
  so $[X,Y]$ is left invariant. Consequently $[X_e,Y_e]=[X,Y]_e$ defines the Lie
  algebra $\mathcal{G}=T_eG$. Formula (2) is left invariant by construction.
  For $a\in G$, the inner automorphism $R_{a^{-1}}L_a$ fixes $e$ and has
  differential $\operatorname{Ad}(a)=dR_{a^{-1}}dL_a$. If $x_t$ is the flow of a
  left-invariant $X$, then $L_y\circ x_t=x_t\circ L_y$, hence
  $x_t(y)=R_{x_t(e)}(y)$ and $dx_t=dR_{x_t(e)}$. Using
  $[Y,X]=\lim_{t\to0}t^{-1}(dx_t(Y)-Y)$ from
  \cref{prop:dc-ch0-5-4} and
  $\operatorname{Ad}(a)Y=dR_{a^{-1}}Y$ gives
  $$
  [Y,X]=\lim_{t\to0}\frac{1}{t}
  \left(\operatorname{Ad}((x_t(e))^{-1})Y-Y\right).
  $$
  Bi-invariance gives
  $\langle U,V\rangle=\langle dR_{x_t(e)}U,dR_{x_t(e)}V\rangle$; differentiating at
  $t=0$ yields (3).
\end{proof}

\begin{proposition}[existence of a left-invariant metric]
  \label{prop:dc-ch1-2-6-leftinv}
  \lean{Riemannian.DCLeftInvariantForm, Riemannian.DCLeftInvariantMetric}
  \uses{def:dc-ch1-2-1}
  \leanok
  Let $G$ be a Lie group with Lie algebra $\mathcal{G}=T_eG$, and let
  $\langle\ ,\ \rangle_e$ be a positive-definite symmetric inner product on
  $\mathcal{G}$. Then (2) defines a left-invariant Riemannian metric on $G$.
\end{proposition}
\begin{proof}
  \leanok
  \uses{def:dc-ch1-2-1}
  \sketch
  Each $(dL_{x^{-1}})_x$ is a linear isomorphism, since $L_x$ and $L_{x^{-1}}$
  are smooth inverses. Pulling back the fixed inner product therefore gives a
  symmetric, bilinear, positive-definite form. Smoothness follows from the smooth
  map $(x,y)\mapsto x^{-1}y$ and its parametrized differential. Formula (2) is
  left invariant by construction.
\end{proof}

\begin{example}[the product metric]
  \label{ex:dc-ch1-2-7}
  \lean{Riemannian.DCProductForm, Riemannian.DCProductMetric}
  \uses{def:dc-ch1-2-1, ex:dc-ch1-2-5}
  \leanok
  \dcref{ch1:2.7}
  For Riemannian manifolds $M_1,M_2$ with projections $\pi_1,\pi_2$, define
  the \emph{product metric} on $M_1\times M_2$ by
  $$
  \langle u,v\rangle_{(p,q)}=
  \langle d\pi_1\cdot u,d\pi_1\cdot v\rangle_p+
  \langle d\pi_2\cdot u,d\pi_2\cdot v\rangle_q.
  $$
  This holds for $(p,q)\in M_1\times M_2$ and
  $u,v\in T_{(p,q)}(M_1\times M_2)$.
  The product of the induced metrics on $S^1\times\cdots\times S^1=T^n$ is the
  \emph{flat torus} metric.
\end{example}
\begin{proof}
  \leanok
  \uses{ex:dc-ch1-2-5}
  \sketch
  Each summand is the pullback of a factor metric along a projection and is
  symmetric, bilinear, and differentiable by \cref{ex:dc-ch1-2-5}; their sum has
  the same properties. If $u\ne0$, its two projected components are not both zero,
  so positive definiteness of one factor makes the sum positive. Thus
  \cref{def:dc-ch1-2-1} applies.
\end{proof}

\begin{definition}[parametrized curve]
  \label{def:dc-ch1-2-8}
  \lean{Riemannian.DCIsCurve}
  \dcref{ch1:2.8}
  A differentiable map $c:I\to M$ from an open interval $I\subset\mathbb{R}$ to
  a differentiable manifold is a \emph{(parametrized) curve}; it may have
  self-intersections or corners.
\end{definition}

\begin{definition}[vector field along a curve; length]
  \label{def:dc-ch1-2-9}
  \lean{Riemannian.DCIsVectorFieldAlong, Riemannian.DCVelocity, Riemannian.DCArcLength}
  \uses{def:dc-ch1-2-1, def:dc-ch1-2-8}
  \leanok
  \dcref{ch1:2.9}
  A \emph{vector field along a curve} $c:I\to M$ is a differentiable assignment
  $V(t)\in T_{c(t)}M$, where differentiability means that $t\mapsto V(t)f$ is
  differentiable for every differentiable $f$ on $M$. The velocity field is
  $dc(\frac{d}{dt})=\frac{dc}{dt}$. For $[a,b]\subset I$, the restriction is a
  \emph{segment} and its length is
  $$
  \ell_a^b(c)=\int_a^b\left\langle\frac{dc}{dt},\frac{dc}{dt}\right\rangle^{1/2}\,dt.
  $$
\end{definition}

\begin{lemma}[isometries preserve length]
  \label{lem:dc-ch1-2-9-isom}
  \lean{Riemannian.DCIsometry.dcArcLength, Riemannian.DCPreservesMetric.dcArcLength}
  \uses{def:dc-ch1-2-2, def:dc-ch1-2-9}
  \leanok
  If $f:M\to N$ is an isometry and $c:I\to M$ is differentiable, then
  $$
  \ell_a^b(f\circ c)=\ell_a^b(c)
  $$
  for every segment $[a,b]\subset I$.
\end{lemma}
\begin{proof}
  \leanok
  \sketch
  The chain rule gives $\frac{d(f\circ c)}{dt}=df_{c(t)}(\frac{dc}{dt})$. Applying
  (1) to the velocity vector makes the two length integrands equal pointwise, so
  their integrals agree.
\end{proof}

\begin{proposition}[existence of Riemannian metrics]
  \label{prop:dc-ch1-2-10}
  \lean{Riemannian.exists_riemannianMetric}
  \uses{def:dc-ch1-2-1}
  \leanok
  \dcref{ch1:2.10}
  Every differentiable Hausdorff manifold with a countable basis admits a
  Riemannian metric.
\end{proposition}
\begin{proof}
  \leanok
  \uses{thm:dc-ch0-5-6,ex:dc-ch1-2-4}
  \sketch
  By \cref{thm:dc-ch0-5-6}, choose a differentiable partition of unity
  $\{f_\alpha\}$ subordinate to a locally finite cover by coordinate neighborhoods
  $V_\alpha$, with $f_\alpha\ge0$, $\mathrm{supp}\,f_\alpha\subset V_\alpha$,
  and $\sum_\alpha f_\alpha=1$. In a chart
  $\mathbf{x}_\alpha:V_\alpha\to\mathbb{R}^n$, pull back the Euclidean metric to
  obtain
  $$
  \langle u,v\rangle_p^\alpha=
  \left\langle d(\mathbf{x}_\alpha)_p u,d(\mathbf{x}_\alpha)_p v\right\rangle,
  \qquad p\in V_\alpha.
  $$
  Each local form is symmetric, bilinear, positive definite, and differentiable.
  Define
  $$
  \langle u,v\rangle_p=\sum_\alpha f_\alpha(p)\langle u,v\rangle_p^\alpha,
  $$
  omitting terms with $f_\alpha(p)=0$. Local finiteness gives differentiability and
  a finite sum near each point. Nonnegativity of every summand and
  $\sum_\alpha f_\alpha(p)=1$ imply positivity for $u\ne0$ from any index with
  $f_\alpha(p)>0$. Thus this is a Riemannian metric.
\end{proof}

\subsection{Volume on an oriented manifold}

For a positive parametrization $\mathbf{x}:U\to M$ and a positive orthonormal
basis of $T_pM$, write
$X_i(p)=\frac{\partial}{\partial x_i}(p)=\sum_j a_{ij}e_j$. Then
$g_{ik}(p)=\sum_j a_{ij}a_{kj}$ and
$$
\mathrm{vol}(X_1(p),\ldots,X_n(p))=\det(a_{ij})=\sqrt{\det(g_{ij})}(p).
$$
For another positive parametrization $\mathbf{y}$ with coefficients $h_{ij}$,
$$
\sqrt{\det(g_{ij})}(p)=J\sqrt{\det(h_{ij})}(p),\qquad
J=\det\left(\frac{\partial y_i}{\partial x_j}\right)>0.\qquad\text{(4)}
$$
If $R\subset M$ is open, connected, has compact closure, lies in a positively
parametrized neighborhood $\mathbf{x}(U)$, and the boundary of
$\mathbf{x}^{-1}(R)$ has measure zero, define
$$
\mathrm{vol}(R)=\int_{\mathbf{x}^{-1}(R)}\sqrt{\det(g_{ij})}\,dx_1\cdots dx_n.\qquad\text{(5)}
$$
The change-of-variables theorem and (4) make this independent of the chosen
positive chart; orientability fixes the sign.

\begin{remark}[volume form]
  \label{rem:dc-ch1-2-11}
  \uses{lem:dc-ch1-2-11-pointwise}
  \dcref{ch1:2.11}
  Equation (4) identifies the integrand in (5) with a positive differential
  $n$-form, the \emph{volume form} (or volume element) $\nu$. For a compact region
  not contained in one chart, choose a partition of unity $\{\varphi_i\}$ subordinate
  to a finite chart cover and set
  $$
  \mathrm{vol}(R)=\sum_i\int_{\mathbf{x}_i^{-1}(R)}\varphi_i\nu;
  $$
  the result is independent of the partition.
\end{remark}

\begin{remark}[volume from a volume element]
  \label{rem:dc-ch1-2-12}
  \dcref{ch1:2.12}
  A globally defined positive differential $n$-form gives a notion of volume on a
  differentiable manifold; a Riemannian metric is only one way to obtain such a
  volume element.
\end{remark}
